\section{Data, Training, and Evaluation}
\label{sec:experiments}

\subsection{Training data and split boundary}

We use PLINDER 2024-06/v2 \citep{durairaj2024plinder}. The preserved processed
pool contains 50,213 complexes. The current compatibility split contains
47,310 training and 1,076 validation sample keys. Exact sample-key, canonical
SMILES, and PLINDER pocket70-community overlaps between these two partitions
are zero. Ligands contain between 5 and 120 heavy atoms and at most 30
fragments; cropped receptors contain at least 50 residues.

\paragraph{Draft scientific boundary.}
The preserved split predates the final strict exclusion of canonical ligand
SMILES from both external benchmarks. We therefore describe the current
external results as frozen compatibility diagnostics, not as evidence of
prospective or chemically novel generalization. A publishable retraining claim
requires the strict split manifest and leakage audit described in the project
data contract.

\subsection{Optimization}

The current geometry training recipe uses six equivariant layers, scalar width
384, 32 channels for each \(l=1\) parity, 16 channels for each \(l=2\)
parity, 32 radial basis functions, dropout 0.1, and a \(5\,\angstrom\)
dynamic-contact cutoff. We use AdamW with learning rate \(3\times10^{-4}\),
weight decay 0.01, gradient-norm clipping at 1.0, and a warmup--stable--cosine
schedule over at most 200,000 updates. The per-rank batch size is 64,
size-aware batching balances variable graph sizes, and exponential moving
average weights use decay 0.999. Automatic mixed precision is disabled for the
current fused equivariant kernels; TF32 is enabled.

The promoted geometry checkpoint contains approximately \(5.72\) million
learned parameters, excluding serialized equivariant-operator constants. The
separate confidence checkpoint contains approximately \(9.40\) million
parameters, for a combined inference stack of approximately \(15.12\) million
parameters.

The confidence model is trained on 4,096 PLINDER complexes with pose sets
generated under the matched \(M=80\), 25-step, \(\sigma=0.5\,\angstrom\),
\(10\,\angstrom\)-pocket protocol. Training uses at most 80 poses per
complex, four confidence interaction layers, hidden width 512, and the loss in
Eq.~\eqref{eq:confidence_loss}. External benchmark labels are not used for
confidence training, checkpoint selection, or hyperparameter selection. The
retained checkpoint is step 42,500; subsequent registered fine-tuning
experiments did not pass their PLINDER validation gate.

\subsection{Frozen reference-pocket redocking}

We evaluate on Astex Diverse (85 complexes)
\citep{hartshorn2007astex} and PoseBusters v2 (308 complexes)
\citep{buttenschoen2024posebusters}. In this diagnostic, the pocket center is derived
from receptor residue virtual nodes lying within \(8\,\angstrom\) of the
reference ligand. The reference ligand is removed before sampling and its
coordinates are subsequently used only to compute the evaluation label. The
model receives canonical ligand chemistry, the cleaned receptor, and the
frozen pocket center. Because the site is target-derived, these numbers do not
measure blind pocket finding.

For every complex, we sample the same 80 poses with 25 ODE steps,
\(\sigma=0.5\,\angstrom\), a \(10\,\angstrom\) pocket crop, no refinement,
and base seed 42 with deterministic complex-ID offsets. We compare:
(i) the first generated pose; (ii) Vina plus the ligand distance-geometry
strain term \citep{trott2010vina}; (iii) pure learned confidence;
(iv) the frozen historical composite selector; and (v) the minimum-RMSD
oracle among the same 80 poses. The oracle uses the reference label and is
reported only as a sampling-coverage upper bound.

\begin{figure}[t]
    \centering
    \includegraphics[width=\textwidth]{protein_pocket_context.png}
    \caption{Explicit-pocket inference boundary illustrated on Astex 1T46.
    The supplied center defines a residue-aware \(10\,\angstrom\) receptor
    crop. The reference ligand is shown only to explain this diagnostic; it is
    removed before sampling. Because the center is target-derived, the
    experiment measures redocking robustness rather than pocket discovery.}
    \label{fig:pocket_context}
\end{figure}

\subsection{Pocket crop and center perturbation}

We characterize the promoted model under the pre-registered V2
pocket-sensitivity protocol. The docking and
confidence checkpoints, 80-candidate budget, 25-step integrator,
\(\sigma=0.5\,\angstrom\), and per-complex random stream are fixed. We vary
the receptor crop radius over \(\{6,8,10,12\}\,\angstrom\) and perturb each
reference-defined center by isotropic Gaussian noise with standard deviation
\(\{0,1,2\}\,\angstrom\). The 1 and \(2\,\angstrom\) perturbations for one
complex share a direction, enabling within-complex interpretation of the
perturbation magnitude. All 24 dataset--condition runs cover all 85 Astex and
308 PoseBusters complexes without unresolved failures. This study
characterizes robustness only; external benchmark outcomes are not used to
select a new default crop radius.

\subsection{Metrics}

The primary pose metric is top-1 symmetry-aware heavy-atom RMSD success below
\(2\,\angstrom\). We additionally report median RMSD, oracle top-80 success,
failure rate, and official PoseBusters validity excluding the RMSD criterion.
Sampling coverage and top-1 selection are reported separately: an oracle
success is not a deployable docking success. Likewise, a docking score is used
for within-candidate selection and is not interpreted as binding affinity.
For primary complex-level success proportions we report 95\% Wilson intervals
where space permits. These intervals describe finite benchmark-set
uncertainty; they do not replace variability across training or sampling
seeds, which has not yet been measured for the promoted stack.
